\section*{Chapter \ref{ch:iv2:valuation-and-financial-statements} --- Valuation and Financial Statements}
\markboth{\protect\omTocMark{}{Answers: Valuation and Financial Statements}}{\protect\omTocMark{}{Answers: Valuation and Financial Statements}}

\begin{solution}{iq:iv2:valuation-and-financial-statements:1}
Net income falls by 7.5. Operating cash flow and cash rise by 2.5. Fixed assets fall by 10, and retained earnings fall by 7.5.

Depreciation is an expense, so pre-tax profit falls by 10, and tax falls by 2.5, which leaves net income down 7.5. On the cash flow statement you start from net income and add back the extra depreciation, because it is not cash, so $-7.5 + 10 = +2.5$: the only real cash effect is the tax saved. On the balance sheet, assets change by $+2.5 - 10 = -7.5$, and equity changes by $-7.5$ through retained earnings. So it balances.

\iqwrong{``Depreciation is non-cash, so cash does not change.'' The expense itself is not cash, but it lowers the tax you pay, and that is.}
\iqfollow{What if it were 10 of extra interest, paid in cash?}{Net income still falls by 7.5, and there is nothing to add back, so cash also falls by 7.5. Under US GAAP interest paid is in operating cash flow, so that falls by 7.5 too. IFRS has let a company show it under financing instead, so operating cash flow need not fall, and IFRS 18, effective from 2027, makes financing the rule for most companies. The liabilities do not change, and equity falls by 7.5.}
\end{solution}

\begin{solution}{iq:iv2:valuation-and-financial-statements:2}
Net income is 30, and operating cash flow is $-10$. Profit is positive while cash falls.

Pre-tax profit is $100 - 60 = 40$, tax is 10, and net income is 30. For cash flow, start at 30. Receivables rise by 100, which you subtract. Inventory rises and falls by 60, so it nets to zero. Payables rise by 60, which you add. So $30 - 100 + 60 = -10$, which is the tax paid. Sales and purchases have moved no cash.

Assets are receivables of 100 and cash of $-10$, so 90. Liabilities are payables of 60, and equity is retained earnings of 30. And $90 = 60 + 30$.

\iqwrong{Equating profit with cash. A company that grows on credit can report rising profits and run out of money, and many do.}
\end{solution}

\begin{solution}{iq:iv2:valuation-and-financial-statements:3}
Enterprise value is $1000 + 300 - 100 + 50 + 30 = 1{,}280$, and the multiple is 8.0 times. None of the three transactions changes the enterprise value.

Enterprise value\index{enterprise value} is the value of the operations to everyone who has a claim on them. It adds the other claims to equity and takes off the cash, which the operations do not need. It matches EBITDA, a number before any payment to capital providers.

If it borrows 100, then debt and cash both rise by 100. If it pays a dividend of 50, market capitalisation and cash both fall by 50. If it buys back 100 of shares, market capitalisation and cash fall by 100. In each case operations have not changed, and that is the point of the measure. In the last two cases the market capitalisation falls because value has left the company. It does not fall because of bad news.

\iqwrong{Mixing up claim-holders: comparing enterprise value with net income, or a share price with EBITDA. A multiple's numerator and denominator must belong to the same claim-holders.}
\end{solution}

\begin{solution}{iq:iv2:valuation-and-financial-statements:4}
Free cash flow to the firm is 100. Free cash flow to equity is 95. You discount the first at the weighted average cost of capital, because it belongs to all the capital providers.

To the firm: tax operating profit as if there were no debt, $200 \times 0.75 = 150$, add back the non-cash 40, and take off the 70 of investment and the 20 of working capital. That gives 100.

To equity: take off the interest after its tax shield, $20 \times 0.75 = 15$, and add the 10 of new borrowing, which gives 95. A check from the top: net income is $(200 - 20) \times 0.75 = 135$, and $135 + 40 - 70 - 20 + 10 = 95$.

\iqwrong{Deducting interest and then discounting at the weighted average cost of capital. The rate already includes the cost of debt and its tax shield, so you would count interest twice.}
\iqfollow{And free cash flow to equity?}{Discount it at the cost of equity, which gives equity value directly. Done with consistent assumptions, the two routes agree.}
\end{solution}

\begin{solution}{iq:iv2:valuation-and-financial-statements:5}
It is worth $100/(0.09 - 0.03) \approx 1{,}667$. With 4\% growth, or with an 8\% cost of capital, it is worth 2,000, which is 20\% more. If 100 is this year's flow, next year's is 103, and the value is about 1,717.

This is the Gordon growth\index{Gordon growth model} formula: value is next year's flow divided by the spread between the discount rate and growth. The spread here is 6 points, so moving it by one point is a large change. A 2\% growth rate gives $100/0.07 \approx 1{,}429$, a fall of 14\%. That is smaller than the 20\% rise from 4\% growth, because the value is convex in the spread.

\iqwrong{Using this year's flow in the numerator. The formula discounts a flow arriving one year from today, and using this year's understates value by the growth rate.}
\iqfollow{What multiple does this imply?}{About 16.7 times next year's free cash flow. Any multiple hides a spread between a discount rate and a growth rate, and it is worth asking what they are.}
\end{solution}

\begin{solution}{iq:iv2:valuation-and-financial-statements:6}
There are 118 diluted shares, and the diluted equity value is 2,950. Because the convertible is treated as shares, you take its 200 out of debt in the bridge to enterprise value.

Options, by the treasury stock method\index{treasury stock method}: if exercised, they bring in $10 \times 10 = 100$, which buys back 4 shares at 25, for a net addition of 6. Restricted shares add 2 in full, as there is nothing to pay. The convertible is in the money, because 25 is above 20, so it becomes $200/20 = 10$ shares. So $100 + 6 + 2 + 10 = 118$, and at 25 that is 2,950, against 2,500 undiluted.

\iqwrong{Counting the 200 twice, as 10 shares in equity value and again as debt. Choose one, according to whether the convertible is in the money.}
\iqfollow{And at a price of 15?}{The convertible stays as debt of 200, and the options add only $10 - 100/15 \approx 3.33$ shares. The count depends on the price, and in a tight valuation that makes the calculation circular.}
\end{solution}

\begin{solution}{iq:iv2:valuation-and-financial-statements:7}
Operating profit rises by 40\%, from 20 to 28. Pre-tax profit rises by 80\%, from 10 to 18. Operating profit reaches zero at revenue of 75, and pre-tax profit at revenue of 87.5.

Each extra unit of revenue adds 0.8 of profit, so 10 of extra revenue adds 8. The base of 20 is small, because fixed costs take most of the contribution. Operating leverage\index{operating leverage} is the ratio of contribution to operating profit, $80/20 = 4$. Financial leverage is $20/10 = 2$. Together they multiply to 8, so a 10\% rise in revenue gives 80\%.

For break-even, $0.8R = 60$ gives $R = 75$, and $0.8R = 70$ gives $R = 87.5$. Revenue only has to fall 12.5\% to wipe out the pre-tax profit.

\iqwrong{Assuming leverage works in one direction. The multiplier works the same on the way down. A 10\% fall in revenue takes pre-tax profit from 10 to 2.}
\end{solution}

\begin{solution}{iq:iv2:valuation-and-financial-statements:8}
The cost of equity is 10\%, and the weighted average cost of capital is 8.75\%. With equal debt and equity, the equity beta is about 1.77, the cost of equity is about 12.84\%, and the weighted average cost of capital is about 8.67\%.

By the capital asset pricing model\index{capital asset pricing model}, the cost of equity is $4\% + 1.2 \times 5\% = 10\%$. After tax, debt costs $5\% \times 0.75 = 3.75\%$. The weights are 80\% and 20\%, so $8\% + 0.75\% = 8.75\%$.

To remove the leverage, $\beta_u = \beta / \left(1 + (1 - t)D/E\right) = 1.2/(1 + 0.75 \times 0.25) \approx 1.01$. To restore it at $D/E = 1$, multiply by $1 + 0.75 = 1.75$, which gives about 1.77. The cost of equity is $4\% + 1.77 \times 5\%$, about 12.84\%. The new weighted average cost of capital is $0.5 \times 12.84\% + 0.5 \times 6\% \times 0.75 \approx 8.67\%$.

\iqwrong{Keeping the cost of equity at 10\% after adding debt. Doing so would make the cost of capital fall to 7.25\%, a free lunch. Debt makes equity riskier, and that gives back most of the gain.}
\end{solution}

\begin{solution}{iq:iv2:valuation-and-financial-statements:9}
The company is worth 1,125, its equity 625, its cost of equity 13\%, and its weighted average cost of capital about 8.89\%. It is not about 7\% because equity gets more expensive as debt rises.

With no debt the company is worth $100/10\% = 1{,}000$. Interest gives a tax saving of $25\% \times 5\% \times 500$ every year forever. Discounted at the debt's own 5\%, that is $tD = 125$. This is the tax shield\index{tax shield}, and the value is $1{,}000 + 125$, as Modigliani and Miller showed with taxes\index{Modigliani-Miller theorem}. Equity is $1{,}125 - 500 = 625$.

Equity holders get $100 - 500 \times 5\% \times 0.75 = 81.25$, so their return is $81.25/625 = 13\%$. Then $\tfrac{625}{1125} \times 13\% + \tfrac{500}{1125} \times 3.75\% \approx 8.89\%$. And $100/8.89\% = 1{,}125$, which closes the loop.

\iqwrong{Holding the cost of equity at 10\%, which gives about 7.22\%. That would make the company worth about 1,385. The tax shield is worth only 125, and debt does not lower the risk of the operations.}
\iqfollow{So should it borrow without limit?}{Not in reality. The stem removes the cost of distress, and in practice this cost, and a rising cost of debt, set the limit.}
\end{solution}

\begin{solution}{iq:iv2:valuation-and-financial-statements:10}
It is worth about 2,101, and the terminal value makes up about 78\%. At 4\% terminal growth it is worth about 2,446. If cash arrives in mid-year, the value rises by a factor of $1.09^{1/2} \approx 1.044$, to about 2,193.
\begin{center}\small
\begin{tabular}{@{}lrrrrr@{}}
Year & 1 & 2 & 3 & 4 & 5 \\ \hline
Free cash flow & 100.0 & 110.0 & 121.0 & 133.1 & 146.4 \\
Present value at 9\% & 91.7 & 92.6 & 93.4 & 94.3 & 95.2
\end{tabular}
\end{center}
The five years give about 467. The terminal value at year 5 is $146.41 \times 1.03/(0.09 - 0.03) \approx 2{,}513$. Discounted by $1.09^5 \approx 1.539$, it is about 1,634.

Those present values are nearly equal because 10\% growth almost cancels 9\% discounting. So a quick sum is about $5 \times 93$.

\iqwrong{Saying that a 78\% terminal share is a flaw in the method. It is what long-lived businesses look like. It means your attention belongs on the assumptions behind that number, the long-run margin and the return on reinvestment, and not on the third decimal of year 2.}
\iqfollow{One more point of terminal growth lifts the terminal value from about 2,513 to about 3,045. Is that plausible?}{That is 21\% more, for one point of growth. Check it against the reinvestment that growth requires, or you are getting growth for free.}
\end{solution}

\begin{solution}{iq:iv2:valuation-and-financial-statements:11}
A multiple of 10 implies growth of about 3.32\%. Growth of 2\% implies about 8.0 times. A multiple of 18 requires about 5.8\%, which is too high to last forever.

The terminal value is $\mathrm{FCF}(1 + g)/(r - g)$. Setting it equal to $TV$ and solving gives
\[
g = \frac{TV \cdot r - \mathrm{FCF}}{TV + \mathrm{FCF}} = \frac{1000 \times 0.09 - 55}{1000 + 55} \approx 3.32\%.
\]
At $g = 2\%$, the terminal value is $55 \times 1.02/0.07 \approx 801$, about 8.0 times EBITDA. At 1,800, the same formula gives $(162 - 55)/1855 \approx 5.8\%$.

A company that grows at 5.8\% forever would eventually outgrow any economy. So that growth is what you say: 18 times assumes a rate above long-run nominal growth, unless margins or capital needs change too.

\iqwrong{Treating the multiple and the growth rate as two independent checks. They are one assumption written twice, so always convert one into the other, and look at what it asks you to believe.}
\iqfollow{Is the 55 fair as the base?}{Ask how much of it is depressed by growth investment in year 5. A company slowing down needs less reinvestment, and its base flow rises, so the implied growth falls.}
\end{solution}

\begin{solution}{iq:iv2:valuation-and-financial-statements:12}
It is worth 0 at a 5\% return, 1,000 at 10\%, and 1,500 at 20\%. With no growth, it is worth 1,000. Growth adds value only when the return on new investment is above the cost of capital.

To grow at $g$ you must reinvest a fraction $g/\mathrm{ROIC}$ of profit, because new profit equals investment times the return. So free cash flow is $\mathrm{NOPAT}(1 - g/\mathrm{ROIC})$, and
\[
V = \frac{\mathrm{NOPAT}\left(1 - g/\mathrm{ROIC}\right)}{r - g}.
\]
With a return of 20\%, you reinvest 25 out of 100, which adds $25 \times 20\% = 5$ of profit, or 5\% growth. The cash flow is 75, and $75/5\% = 1{,}500$. With a return of 10\%, the cash flow is 50, and $50/5\% = 1{,}000$, the same as no growth. So growth at the cost of capital is worth nothing. With a return of 5\%, all the profit is reinvested, and no cash ever comes out.

\iqwrong{Using a higher growth rate to justify a higher price without asking what it takes. A DCF with growth and no reinvestment assumes a free lunch.}
\iqfollow{How does this link to the price to earnings multiple?}{With return on equity in place of ROIC, the forward multiple is $(1 - g/\mathrm{ROE})/(r - g)$. For a return of 15\%, a cost of equity of 10\% and growth of 5\%, that is about 13.3.}
\iqwhy{This is the single most useful formula for telling a good business from a large one. It says when growth is worth paying for.}
\end{solution}

\begin{solution}{iq:iv2:valuation-and-financial-statements:13}
EV to EBITDA is 10 times for both. The price to earnings multiple is 16 for $A$ and about 11.4 for $B$. It differs because the two have different leverage. With debt at 12\%, $B$'s multiple is about 28.6.

EBITDA is 120 for both, and $1200/120 = 10$. Company $A$ has net income of $100 \times 0.75 = 75$ and equity of 1,200, so 16 times. Company $B$ pays interest of 30, so its net income is 52.5, and its equity is 600, which is about 11.4 times.

The rule compares the unlevered earnings yield with the after-tax cost of debt. $A$'s earnings yield is $75/1200 = 6.25\%$, and debt after tax costs 3.75\%. Replacing expensive equity by cheaper debt raises earnings per unit of equity, so the multiple falls. At 12\%, debt costs 9\% after tax, which is above 6.25\%. Net income is $(100 - 72) \times 0.75 = 21$, and the multiple rises.

\iqwrong{Calling $B$ cheap because its price to earnings multiple is lower. The two have exactly the same operations and the same enterprise value. $B$'s lower multiple just reflects riskier equity.}
\end{solution}

\begin{solution}{iq:iv2:valuation-and-financial-statements:14}
Free cash flow is 100 either way. Expensing gives EBITDA, operating profit and operating cash flow of 100, and investing cash flow of 0. Capitalising gives EBITDA of 200, operating profit of 180, operating cash flow of 200 and investing cash flow of $-100$. In a steady state operating profit is equal again. With spending growing 20\% a year, it stays above, by about 70 in year 6.

Capitalising moves 100 from an operating cost to an investment, so EBITDA and operating cash flow both improve by the whole 100. Operating profit bears one fifth of it, 20, and improves by 80. Nothing about the cash has changed.

With 100 every year, after five years the amortisation is $5 \times 20 = 100$, the same as expensing, although EBITDA stays 100 higher. With growth of 20\%, year 6 spending is about 248.8 and amortisation is about 178.6, the average of the last five years' spending, and the gap is about 70. A growing capitaliser keeps reporting more than a company that expenses.

Under IFRS, research is expensed, and development is capitalised when it meets specified criteria, so two companies can really differ.

\iqwrong{Comparing the two on EBITDA or operating cash flow, without adjusting. Compare on free cash flow, or put both on one policy.}
\iqfollow{What would raise your suspicion?}{A jump in the capitalised amount, as a share of revenue or of total spending, with no new product to explain it.}
\end{solution}

\begin{solution}{iq:iv2:valuation-and-financial-statements:15}
The cash conversion cycle\index{cash conversion cycle} is 75 days, and net working capital is about 172.6. Growth of 20\% absorbs about 34.5, and cutting 15 days of inventory releases about 24.7.

Receivables are $1000 \times 45/365 \approx 123.3$. Inventory is $600 \times 60/365 \approx 98.6$. Payables are $600 \times 30/365 \approx 49.3$. So net working capital is $123.3 + 98.6 - 49.3 = 172.6$. The cycle is $45 + 60 - 30 = 75$.

With days unchanged, working capital scales with sales, so growth of 20\% needs 20\% more, which is 34.5 of cash. Fifteen days of inventory are $600 \times 15/365 \approx 24.7$.

\iqwrong{Calling growth free of cash. Here each 100 of extra revenue needs about 17 of working capital. Fast growth with a long cycle is a common reason why profitable companies run out of cash.}
\iqfollow{What if the cycle were negative?}{Then customers pay before suppliers are paid, and growth produces cash. Subscription and some retail businesses can work like that, and that cash stops when growth stops.}
\end{solution}

\begin{solution}{iq:iv2:valuation-and-financial-statements:16}
Management quotes a yield of 5\%. Treating share-based pay as a cost, it is about 3.33\%. Keeping the share count flat takes about 1.67 shares a year. At 20 times, the two valuations are 3,000 and 2,000, so the higher is 50\% above the lower.

Market capitalisation is $100 \times 30 = 3{,}000$. The reported yield is $150/3000$, and the adjusted one is $100/3000$. New shares of 50 at 30 are about 1.67 shares, or about 1.67\% dilution a year. Buying them back would cost the 50 in cash, which makes the true free cash flow 100.

IFRS requires share-based pay to be recognised as an expense in the accounts. It is added back in cash flow because no cash leaves. But employees have still been paid, and the holders have paid them, through dilution.

\iqwrong{``It is non-cash, so it is not a real cost.'' You can pay the cost in cash through a buyback, or in dilution by not buying back. To ignore both is to value the company as if its staff worked for free.}
\iqfollow{What if you deduct it and also use a diluted share count?}{Count future grants once only. Existing awards go in the share count, and future awards go in as a cost. Using both for the same award is double counting.}
\end{solution}

\begin{solution}{iq:iv2:valuation-and-financial-statements:17}
With shares, earnings per share are $120/115 \approx 1.043$, up 4.3\%. With debt, they are $106.5/100 = 1.065$, up 6.5\%. At a price of 450 in shares, it falls to about 0.98. A deal that adds to earnings per share does not by that fact create value.

Company $A$ earns 1 per share, at a price to earnings multiple of 20. In shares, it issues $300/20 = 15$ shares. In debt, interest is $300 \times 6\% = 18$, which is 13.5 after tax, so net income is $120 - 13.5 = 106.5$.

The rule for shares is to compare multiples: a target bought at 15 times is accretive\index{accretion and dilution} for a buyer at 20 times. At 450, the price is 22.5 times, above 20, and so it dilutes. The rule for debt is to compare yields: $B$'s earnings yield at this price is $20/300 \approx 6.67\%$, above the after-tax cost of debt of 4.5\%.

\iqwrong{Reading accretion as value. Earnings per share rise whenever you buy a lower multiple, which tells you nothing about whether 300 is a good price. The buyer's multiple may fall to reflect a slower or riskier mix, and the holders are no better off.}
\iqfollow{What decides value, then?}{Whether the price is below what $B$ is worth to you, including synergies, on a discounted basis. Compare the return on the 300 with the cost of capital of the target.}
\end{solution}

\begin{solution}{iq:iv2:valuation-and-financial-statements:18}
Equity at exit is about 976, a multiple of 2.44 on 400, and an annual return of about 19.5\%. At 9 times, equity is about 849, a multiple of 2.12, and the return is about 16.2\%.

EBITDA grows to $100 \times 1.05^5 \approx 127.6$. At 10 times, the enterprise value is about 1,276. Debt has fallen to 300, so equity is about 976. The return is $2.44^{1/5} - 1$. A quick check: at 15\% money doubles in about five years, and 2.44 is a bit more than one doubling, so the return is a bit above 15\%.

The total gain of 576 comes from two places: 276 from EBITDA growth at a constant multiple, and 300 from paying off debt. None comes from a change in the multiple. The debt is why the 28\% rise in enterprise value turns into a 144\% rise in equity.

\iqwrong{Assuming the multiple will rise at exit. A model that needs this is a bet on the market, not on the company. So a proper model holds it flat or lower.}
\iqfollow{What limits the debt?}{The ability to pay interest and repay in a bad year. Check interest cover under a fall in EBITDA, before you look at the return.}
\end{solution}

\begin{solution}{iq:iv2:valuation-and-financial-statements:19}
The lease liability and the asset are both about 421.2. In year 1, depreciation is about 84.2 and interest about 25.3, EBITDA rises by 100, and pre-tax profit falls by about 9.5. Enterprise value including the lease is about 1,421, and EBITDA is 250, which gives about 5.68 times, down from 6.67. Cash does not change. Total expense over 5 years stays at 500.

The present value of five payments of 100 at 6\% is 421.24. The asset is depreciated evenly over 5 years, $421.24/5 = 84.25$. Interest in year 1 is $6\% \times 421.24 = 25.27$. Together they come to 109.52, against rent of 100. Interest falls as the liability falls, to 346.51 after year 1 and on to zero, so the expense is higher early and lower late: from 109.52 in year 1 to 89.91 in year 5.

EBITDA gains the whole rent, because rent has left operating costs. Operating profit gains only $100 - 84.25 = 15.75$. The cash payment of 100 is the same, but of it, $100 - 25.27 = 74.73$ now counts as repaying a debt.

This is IFRS 16\index{lease accounting}. Under the US model, such a lease goes on the balance sheet too, but the expense typically stays a straight-line operating cost.

\iqwrong{Thinking the company looks better. EBITDA is up, but so is debt. If you compare a multiple between companies, both must be on one basis, or the lower multiple is an illusion.}
\end{solution}

\begin{solution}{iq:iv2:valuation-and-financial-statements:20}
Net income is 30 every year. The rest changes by year:
\begin{center}\small
\begin{tabular}{@{}lrrrrr@{}}
Year & 1 & 2 & 3 & 4 & 5 \\ \hline
Net income & 30.0 & 30.0 & 30.0 & 30.0 & 30.0 \\
Cash tax & 2.5 & 2.5 & 15.0 & 15.0 & 15.0 \\
Deferred tax balance & 7.5 & 15.0 & 10.0 & 5.0 & 0.0 \\
Operating cash flow & 57.5 & 57.5 & 45.0 & 45.0 & 45.0
\end{tabular}
\end{center}
At 8\%, paying later instead of evenly is worth about 2.3.

In the accounts, pre-tax profit is $60 - 20 = 40$, tax expense is 10, and net income is 30. For tax in years 1 and 2, the deduction is 50, so taxable profit is 10 and cash tax is 2.5. In years 3 to 5, there is no deduction, so cash tax is 15. The expense of 10 exceeds the cash tax by 7.5 in the first two years, and that builds a liability, the deferred tax\index{deferred tax} balance. It reverses by 5 a year afterwards. Under IAS 12, a company recognises this for temporary differences between the carrying amount of an asset and its tax base.

Total tax paid is 50 either way, the same as $25\% \times 200$. So the gain is only one of timing: $10$ a year discounted, less the actual path, comes to about 2.3.

\iqwrong{Calling the tax saved a permanent saving. It reverses, unless the company keeps buying more assets. A company with growing investment can defer for many years, and that is worth more.}
\end{solution}

\begin{solution}{iq:iv2:valuation-and-financial-statements:21}
Goodwill is 225. The write-up lowers net income by 7.5 a year. An impairment of 100 lowers net income by 100, with no effect on cash or on tax.

Goodwill\index{goodwill} is what is left after assets are marked to fair value. The target's net assets at fair value are book equity of 200, plus the write-up of 100, less a deferred tax liability of $25\% \times 100 = 25$. So they come to 275, and $500 - 275 = 225$. Check: the buyer's assets change by $+200 + 100 + 225 - 500 = +25$, which equals the new liability of 25.

Extra depreciation of 10 a year lowers pre-tax profit by 10. Since it is not deductible, cash tax does not change. Reported tax falls by 2.5 as the deferred tax unwinds, so net income falls by 7.5.

IFRS 3 allocates the price to identifiable assets and liabilities at fair value, and the rest goes to goodwill. IAS 36 tests goodwill for impairment every year, and a loss on goodwill is never reversed. Under both IFRS and US rules for listed companies, goodwill is not amortised.

\iqwrong{Forgetting the deferred tax, and answering 200. The write-up is not deductible, so future tax will be higher than accounts show, and the buyer has taken on that liability.}
\iqfollow{What does an impairment tell you?}{That the deal now looks overpriced, or that the business has worsened since. The cash left when the deal closed. So look at it as evidence on management's record, and not as a new loss.}
\end{solution}

\begin{solution}{iq:iv2:valuation-and-financial-statements:22}
The price implies constant growth of about 5.88\% forever. With 3\% growth after 10 years, it implies about 14.3\% a year for the 10 years. At 3\% growth from today, the company is worth about 1,236.

For constant growth, solve $3000 = 60(1 + g)/(0.08 - g)$, which gives $g = 180/3060 = 1/17 \approx 5.88\%$. At a growth rate of 3\%, the value is $61.8/0.05 = 1{,}236$.

The two-stage version has no closed form that helps, so you solve numerically. At 14.3\% for 10 years, cash flow grows by a factor of about 3.8, to about 228. The discounted terminal value then makes up about 72\% of the price.

This is a reverse DCF\index{reverse DCF}. You do not ask what the company is worth. You ask what you must believe to pay this price. The 1,236 says that about 59\% of the price is paying for growth above 3\%. Then compare 14.3\% with what comparable companies have done, and what the reinvestment would have to earn.

\iqwrong{Treating 5.88\% as a forecast. It is above any plausible long-run growth, so the constant-growth form only tells you the market expects a period of faster growth.}
\iqfollow{And if you believe 10\%?}{Then the shares are too expensive on those assumptions, and you need to find what the market sees. It may be margins, not growth.}
\end{solution}

\enlargethispage{-2\baselineskip}%
\begin{solution}{iq:iv2:valuation-and-financial-statements:23}
Revenue is 213, and operating profit is 20.15, up about 11.9\% from 18. One point of same-store sales moves profit by 0.6, about 3.0\% of 20.15. With flat same-store sales, operating profit is 18.35, up about 1.9\%.

Base year: revenue of 200, store profit of 30, central costs of 12, so operating profit of 18.

Next year: existing stores sell $200 \times 1.03 = 206$, and the extra 6 of revenue adds $6 \times 30\% = 1.8$ of store profit. New stores sell at an annual rate of $2 \times 0.7 = 1.4$ each, for half a year, so $10 \times 1.4 \times \tfrac12 = 7$, and at 5\% that earns 0.35. Store profit is $30 + 1.8 + 0.35 = 32.15$. Less 12, that is 20.15.

Each point of same-store sales is 2 of revenue, and at 30\% flow-through that is 0.6. With flat sales, existing stores earn their old 30, the new stores add 0.35, and operating profit is 18.35. If the new stores did not open either, it would be last year's 18 exactly, as it must be. So of the 2.15 of growth, 1.8 comes from same-store sales, about 84\%, and the new stores add little in their first year.

\iqwrong{Applying the average 15\% margin to the change in same-store sales, which gives only 0.9 and an operating profit of 19.25. The 15\% is after all the costs, and much of them do not move with sales. A change in sales falls through at a much higher rate than the average margin, on the way up and on the way down.}
\iqfollow{What next?}{Ask what the new stores must earn once mature for them to be worth opening, and whether they take sales from existing ones.}
\end{solution}

\begin{solution}{iq:iv2:valuation-and-financial-statements:24}
Earnings rose 30\%, but cash fell and every working-capital line moved the wrong way. Free cash flow went from $+50$ to $-50$. With development spending at Year 1's share of revenue, net income would have been about 95.5, down from 100 instead of up 30\%.

The checks, with 365 days:
\begin{center}\small
\begin{tabular}{@{}lrr@{}}
 & Year 1 & Year 2 \\ \hline
Days of receivables & 43.8 & 63.9 \\
Days of inventory & 60.8 & 86.2 \\
Net income less operating cash flow & $-10$ & $+70$ \\
Free cash flow, after capitalised development & $+50$ & $-50$
\end{tabular}
\end{center}
The accruals\index{earnings quality} of 70 are about 6.2\% of average assets, which is large. Had receivables grown with sales, they would be 144, so 66 of them, or 5.5\% of revenue, are excess. That points to revenue booked early, or to customers who are not paying. Inventory is 50 above its old ratio to cost of goods sold, which points to weak sales or goods building up.

Deferred revenue fell by 20. Billings, meaning revenue plus the change in deferred revenue, are 1,180, below revenue of 1,200. So part of this year's revenue is drawn from the old balance and is not being replaced.

Capitalised development is 5.8\% of revenue against 2\%. At 2\%, it would be 24, so 46 more would have been expensed, about 34.5 after tax. I ignore amortisation, which would be small here.

\textbf{Innocent explanations to test.} A large sale near year end would explain receivables, and the January cash receipts would settle it. A new product launch could justify both inventory and development. A big customer on longer terms is a real change, but it is a cost, not growth.

\textbf{Questions for management.} What share of Q4 sales came in the last month? How have terms changed? What is being capitalised and why now? What has been collected since year end?

\iqwrong{Declaring fraud. These are flags, not proof. Every one has a possible real cause. What you owe an interviewer is a ranking, a number on each and the check that would settle it.}
\iqwhy{A short seller and a long holder both read a table like this, and this is the skill that separates them from someone who reads the headline.}
\end{solution}

\begin{solution}{iq:iv2:valuation-and-financial-statements:25}
The old price implied 33\% recovery, and the new price implies about 18.6\%. You need about 32.1\% for a 15\% expected return. At 35\%, your expected return is about 18.3\%. A fresh position sized by the 1\% rule would be about 4.8\% of capital, against about 3\% originally. You have lost 0.9\% of capital so far. My recommendation is to hold or add only if you can say why the guidance cut does not break the recovery case, and to cut otherwise.

Each price is a probability-weighted\index{implied probability} average of its two cases, so $p = (P - \text{bear})/(\text{bull} - \text{bear})$. So $(90 - 60)/90 = 1/3$, and $(63 - 50)/70 \approx 18.6\%$. Then a 15\% return requires $(63 \times 1.15 - 50)/70 \approx 32.1\%$. At 35\%, the value is $0.35 \times 120 + 0.65 \times 50 = 74.5$, and $74.5/63 - 1 \approx 18.3\%$.

Originally, the downside was $30/90 = 33.3\%$, so a 1\% loss limit meant $1\%/33.3\% = 3\%$ of capital. It has lost $3\% \times 27/90 = 0.9\%$. Now the downside is $13/63 \approx 20.6\%$, so the same rule allows about 4.8\%, and the upside is about 90.5\%.

The arithmetic says buy more, but it depends on your 35\%. You were wrong about one thing already, and your new bear case is only your estimate. Before adding, ask what the cut says about the recovery, whether your bear case is now really the floor, and whether you would buy it today with no history. To hold your nerve for the sake of it is not an argument.

\iqwrong{Anchoring on 90, the price you paid. It is no longer information. The other classic is cutting because the position hurts, without a view on the numbers.}
\iqfollow{What would change your mind?}{Name it in advance: for example, the next quarter's margin below a stated level, or customers lost. If it happens, you sell without another debate.}
\iqwhy{Every fundamental interview ends here: a position, a piece of bad news, and whether you think in probabilities and sizes, or defend what you already own.}
\end{solution}
